\section*{Introduction}

Symmetry is the organizing principle of machine learning for atomistic systems. Networks whose internal features transform equivariantly under the symmetries of three-dimensional space are the default architecture for predicting properties of molecules and crystals, from interatomic potentials~\cite{unke2021spookynet,polat2026cliffordip} to equivariant transformers~\cite{fuchs2020se3transformer}, and the community describes them almost interchangeably as ``E(3)-equivariant''.

How much of that symmetry should be hard-wired, and how much left to be learned from data, is one of the field's live disputes~\cite{creed2026six}. Augmentation closes the data-efficiency gap given enough epochs~\cite{brehmer2024scale}, unconstrained potentials can beat physically constrained ones on accuracy and speed~\cite{bigi2026unconstrained}, and approximate rotation-equivariance carries negligible consequences in realistic bulk simulation~\cite{langer2024broken}, while the opposing evidence that equivariance pays at scale is empirical rather than structural~\cite{ngo2025scaling}. That dispute has been argued almost entirely over one kind of symmetry, rotations, and weighed on a single scale, accuracy. There is precedent for judging such a trade against an exact law rather than a benchmark~\cite{bigi2025forces}.

This paper opens a second such axis. The Euclidean group E(3) contains improper operations: reflections and the spatial inversion $I\colon\, \mathbf{r} \mapsto -\mathbf{r}$. A network can be exactly equivariant to all rotations, SO(3), while carrying no information about parity, and a substantial class of deployed models is built this way, including the eSCN lineage~\cite{passaro2023escn} and EquiformerV2 (EqV2)~\cite{liao2024equiformerv2}. The distinction is invisible at the configuration surface: such models expose a maximum angular degree, not a parity flag. Nor is it an artifact of one architecture. The current generation of universal potentials, eSEN~\cite{fu2025esen} and UMA~\cite{wood2025uma}, is built on the same spherical-channel lineage, and EquiformerV3 (EqV3)~\cite{equiformerv3} declares equivariance to SE(3) in its title, a group containing no improper operation at all. Whether a model respects improper operations is, in practice, an implementation detail that neither model cards nor benchmarks report. On this axis the stakes are not measured in accuracy. A model that is slightly wrong about a rotation makes a slightly wrong prediction. A model that is blind to parity can make predictions the physics forbids outright: errors of kind, not of degree.

The physics supplies an unusually clean instrument for making this precise. By Neumann's principle~\cite{nye1985crystals}, a crystal's macroscopic property tensors are unchanged by each operation of its point group. The piezoelectric tensor $e_{ijk}$, which couples polarization to strain, is a parity-odd rank-3 tensor: under inversion it changes sign. A centrosymmetric crystal, whose space group contains the inversion, therefore demands invariance and sign reversal at once, so the tensor must equal its own negative. The consequence is an algebraic identity, not an empirical trend: the piezoelectric tensor of every centrosymmetric crystal is exactly zero, as for every parity-odd tensor. Since 92 of the 230 space groups are centrosymmetric, this yields a large, label-free test population on which any nonzero prediction is known, without reference data, to be impossible. It converts the abstract question of how much symmetry to build in into a measurement.

Whether an equivariant network can violate the identity is decided by its feature algebra, and the answer is not a matter of degree. A parity-labelled network must predict exactly zero on a centrosymmetric crystal, at any parameter values, random initialization included. Erase the labels and it remains exactly rotation-equivariant, but nothing forbids the impossible output. We prove both statements below.

Three adjacent lines of work make the question timely without answering it. A literature on symmetry breaking starts from Curie's principle~\cite{curie1894symmetry} and develops relaxations permitting less symmetric outputs when the physics demands them~\cite{smidt2021symmetry,kaba2023symmetry}. We probe the complement: for a parity-odd tensor on a centrosymmetric crystal there is no symmetry breaking to model, and a model that escapes the constraint has not gained flexibility but made an error. A complementary line shows that an E(3)-\emph{invariant} model cannot represent chirality~\cite{dumitrescu2024chirality}, the opposite failure of parity bookkeeping. Third, a growing family of crystal-tensor predictors enforces the correct symmetry at the output, through space-group-informed modules, canonicalization, or symmetry-adapted decompositions~\cite{piezotensornet2024}. Whether the features beneath them can represent the constraint is a prior question, one that grows urgent as property heads are attached to pretrained universal potentials~\cite{hung2024dielectric}, whose feature symmetry group is fixed before any downstream target is chosen.

Here we show that this single design decision, the parity labelling of a network's features, governs what deployed models actually predict. The test population is 2,000 centrosymmetric insulators verified with spglib~\cite{togo2024spglib}, on which the true piezoelectric tensor vanishes by symmetry, so any nonzero prediction is a certified error obtained without labels or a reference calculation. A group-theoretic account predicts in advance, with no free parameters, which of these crystals a rotation-only model is nonetheless forced to get right. The measurement is a matched-pair ablation: NequIP~\cite{batzner2022nequip}, Allegro~\cite{musaelian2023allegro} and MACE~\cite{batatia2022mace} are each built in two arms whose sole intended difference is the parity labelling of their features, with unmodified EqV2 as a deployed rotation-only representative. The parity-labelled arms stay below the operating threshold on every crystal; the rotation-only arms predict responses of ordinary magnitude for the large majority, in exact agreement with that account, and at no accuracy advantage where the constraint is inactive. Retraining on the exact answer, up-weighting it and symmetrizing at test time all fail to remove the violations: none of the training strategies tested recovers the constraint where the symmetry is exact. Finally, an audit of the released landscape, spanning the deployed generation, the dedicated crystal-tensor predictors, and heads on frozen pretrained features, finds that most architectures cannot label parity at all, so the guarantee can only be imposed from outside, one quantity at a time. The obstruction belongs to a symmetry class in current use, not to any single model.
